\chapter{Data System and Physical Storage}

本章自底向上探讨数据的数学本质：从比特域出发，经数、字符的解释协议，到线性表的物理组织。数学层只谈集合与映射，位宽、字节、语言类型均为实例化。

\section{最根本元起点}

\begin{itemize}
    \item \textbf{基本域}：$\mathbb{B} = \{0, 1\}$。
    \item \textbf{比特序列空间}：$\mathbb{B}^* = \bigcup_{n=0}^{\infty} \mathbb{B}^n$。所有存储对象均为 $\mathbb{B}^*$ 中的元素。
    \item \textbf{解释协议}：映射 $\mathcal{I}: \mathbb{B}^n \to \mathcal{D}$ 将比特流解释为数学对象 $\mathcal{D}$（整数、字符、地址等）。同一比特串在不同 $\mathcal{I}$ 下可解释为不同对象。
\end{itemize}

\section{数字体系 (Number)}

\subsection{\texorpdfstring{自然数 ($\mathbb{N}$)}{自然数 (N)}}

\begin{itemize}
    \item \textbf{对象}：$\mathbb{N} = \{0, 1, 2, \dots\}$。
    \item \textbf{不变式（定长无符号）}：$n$ 位 $b = b_{n-1}\dots b_0 \in \mathbb{B}^n$，解码
$$\mathcal{I}_{\mathbb{N}}(b) = \sum_{i=0}^{n-1} b_i 2^i, \quad \mathcal{I}_{\mathbb{N}}: \mathbb{B}^n \to [0, 2^n-1]$$
为双射，值域 $[0,2^n-1]$。
    \item \textbf{定理}：加法在环 $\mathbb{Z}_{2^n}$ 中封闭，溢出即模 $2^n$ 回绕。
    \item \textbf{实例化}：网页 \texttt{unsigned-int} 按 $y \gets \sum b_i2^i$ 逐位累加演示；$n=8$ 时 $y \in [0,255]$ 为实例。
\end{itemize}

\subsection{\texorpdfstring{整数 ($\mathbb{Z}$)}{整数 (Z)}}

\begin{itemize}
    \item \textbf{对象}：$\mathbb{Z} = \{\dots, -2, -1, 0, 1, 2, \dots\}$。
    \item \textbf{不变式（补码）}：$b = b_{n-1}\dots b_0 \in \mathbb{B}^n$，
$$\mathcal{I}_{\mathbb{Z}}(b) = -b_{n-1}2^{n-1} + \sum_{i=0}^{n-2} b_i2^i$$
符号位负权消除 $+0/-0$ 二义性。
    \item \textbf{定理}：$\mathcal{I}_{\mathbb{Z}}: \mathbb{B}^n \to [-2^{n-1}, 2^{n-1}-1]$ 为双射；同一比特加法器对 $\mathcal{I}_{\mathbb{N}}$ 与 $\mathcal{I}_{\mathbb{Z}}$ 均正确（模 $2^n$ 同构）。
    \item \textbf{实例化}：网页 \texttt{twos-complement} 演示符号位 $-2^{n-1}$ 与低位累加；$n=8$ 时 $[-128,127]$ 为实例。
\end{itemize}

\subsection{\texorpdfstring{实数/有理数 ($\mathbb{R} \to \mathbb{Q}_{\text{discrete}}$)}{实数/有理数 (R 到 Q discrete)}}

\textbf{物理事实}：离散状态机无法精确表示连续无理数，实数离散化为 $\mathbb{Q}_{\text{discrete}} \subset \mathbb{Q}$。

\begin{enumerate}
    \item \textbf{定点数 (Fixed-Point)}：$x = m \cdot B^{-f}$，$m \in \mathbb{Z}$，$f$ 固定。值域与精度由 $f$ 权衡。
    \item \textbf{浮点数 (Floating-Point)}：
    \begin{itemize}
        \item \textbf{通用形式}：$x = (-1)^s \cdot m \cdot 2^e$，$s \in \mathbb{B}$，$m$ 为尾数，$e$ 为指数。
        \item \textbf{IEEE 754 结构}：$n$ 位划分为 $[s \mid E \mid F]$，$k$ 为指数位宽，$p$ 为尾数位宽，偏置 $\text{Bias}=2^{k-1}-1$，$E_{\text{raw}} \in [0,2^k-1]$。
        \begin{itemize}
            \item $E_{\text{raw}}=0, F=0$：$\pm 0$。
            \item $E_{\text{raw}}=0, F\neq0$：非规格化 $x=(-1)^s\cdot (F/2^p)\cdot 2^{1-\text{Bias}}$。
            \item $1\le E_{\text{raw}}\le 2^k-2$：规格化 $x=(-1)^s\cdot (1+F/2^p)\cdot 2^{E_{\text{raw}}-\text{Bias}}$（隐含前导 $1$）。
            \item $E_{\text{raw}}=2^k-1, F=0$：$\pm\infty$；$F\neq0$：$\text{NaN}$。
        \end{itemize}
        \item \textbf{实例参数}：binary32 $k=8,p=23,\text{Bias}=127$；binary64 $k=11,p=52,\text{Bias}=1023$。
        \item \textbf{等价判定}：引入容差 $\epsilon>0$，$a\equiv_\epsilon b \iff |a-b|<\epsilon$，避免直接判等。
    \end{itemize}
\end{enumerate}

\begin{itemize}
    \item \textbf{实例化}：网页 \texttt{ieee754} 支持 binary32/64 编解码与分类演示；上述公式为数学层，位串布局为实例。
\end{itemize}

\section{字符体系 (Character)}

\begin{itemize}
    \item \textbf{对象}：有限字母表 $\Sigma$，字符串幺半群 $(\Sigma^*,\circ,\varepsilon)$。
\end{itemize}

\subsection{ANSI / ASCII}

\begin{itemize}
    \item \textbf{对象}：映射 $\phi: \{0,\dots,127\} \to \Sigma_{\text{ASCII}}$ 为双射，$|\Sigma_{\text{ASCII}}|=128$。
    \item \textbf{不变式}：$7$ 位码 $\in [0,127]$。
    \item \textbf{实例化}：存储实例将 $7$ 位以 $0xxxxxxx$ 置于 $1$ 字节（$8$ 位），最高位恒 $0$；网页 \texttt{character-encoding} 模式 \texttt{ascii} 演示该实例。
\end{itemize}

\subsection{Unicode 架构}
\begin{enumerate}
    \item \textbf{抽象字符集 (UCS)}：全局空间 $\mathcal{U}$，$c\in\mathcal{U}$ 对应唯一码点 $u\in\mathbb{N}$，记 $\text{U}+XXXX$。
    \item \textbf{编码方案 (UTF)}：$\psi: \mathcal{U}\to\mathbb{B}^*$ 具前缀可解码性。
    \begin{itemize}
        \item \textbf{UTF-8}：变长 $1$--$4$ 字节，
        \begin{center}
        \small
        $U+0000$--$007F$: $0xxxxxxx$; $U+0080$--$07FF$: $110xxxxx\;10xxxxxx$; $U+0800$--$FFFF$: $1110xxxx\;10xxxxxx\;10xxxxxx$; $U+10000$--$10FFFF$: $11110xxx\;10xxxxxx\;10xxxxxx\;10xxxxxx$
        \end{center}
        向下兼容 ASCII（首字节 $0$ 判单字节）。
        \item \textbf{UTF-16/32}：$\psi$ 的另两种实例化，码点到码元的定/变长映射。
    \end{itemize}
\end{enumerate}

\begin{itemize}
    \item \textbf{实例化}：网页 \texttt{character-encoding} 模式 \texttt{utf8} 按码点区间分步生成字节；字符字形与字节存储均为实例。
\end{itemize}

\section{物理存储结构：线性表（Linear List）}

线性表是有限序列 $L = \langle d_1, d_2, \dots, d_n \rangle$，$d_i \in \mathcal{D}$。

\subsection{顺序存储结构（顺序表 / 数组）}

\begin{itemize}
    \item \textbf{对象}：基址 $\beta$，元大小 $\sigma$，连续区间 $[\beta, \beta+n\sigma)$。
    \item \textbf{不变式（随机访问）}：$\text{Addr}(i)=\beta+i\sigma$，$0\le i<n$，时间 $O(1)$。
    \item \textbf{多维压平（行优先）}：$M\times N$ 矩阵逻辑 $(i,j)$ 映为 $\text{Addr}(i,j)=\beta+(iN+j)\sigma$；列优先为对偶实例。
\end{itemize}

\begin{figure}[H]
\centering
\begin{tikzpicture}[
    cell/.style={draw, minimum width=1.4cm, minimum height=0.8cm, font=\ttfamily},
    addr/.style={font=\small},
    arrow/.style={->, >=latex}
]
    \node[cell] (c0) at (0,0) {$d_0$};
    \node[cell] (c1) at (2.6,0) {$d_1$};
    \node[cell] (c2) at (5.2,0) {$d_2$};
    \node[cell] (c3) at (7.8,0) {$d_3$};
    \node[cell] (c4) at (10.4,0) {$\dots$};
    \draw[arrow] (c0.south) -- ++(0,-0.5) node[addr, below] {$\beta$};
    \draw[arrow] (c1.south) -- ++(0,-0.5) node[addr, below] {$\beta+\sigma$};
    \draw[arrow] (c2.south) -- ++(0,-0.5) node[addr, below] {$\beta+2\sigma$};
    \node[addr, anchor=west] at ($(c4.east)+(0.3,0)$) {$\leftarrow$ 物理地址连续，支持 $O(1)$ 随机访问};
\end{tikzpicture}
\caption{一维顺序表物理内存布局}
\label{fig:array_1d}
\end{figure}

\begin{figure}[H]
\centering
\begin{tikzpicture}[
    cell/.style={draw, minimum size=0.75cm, font=\ttfamily},
    bigcell/.style={draw, minimum width=1.1cm, minimum height=0.75cm, font=\ttfamily},
    arrow/.style={->, >=latex, thick},
    brace/.style={decorate, decoration={brace, amplitude=5pt, mirror}}
]
    \matrix (mlog) [matrix of nodes, nodes={cell}, column sep=-\pgflinewidth, row sep=-\pgflinewidth] {
        (0,0) & (0,1) & (0,2) \\
        (1,0) & (1,1) & (1,2) \\
    };
    \node[left=0.3cm of mlog, font=\bfseries] {逻辑结构};

    \draw[arrow] (mlog.south) -- ++(0,-0.7) node[midway, right, font=\small] {行优先压平};

    \matrix (mphy) [matrix of nodes, nodes={bigcell}, column sep=-\pgflinewidth, below=1.1cm of mlog.south] {
        (0,0) & (0,1) & (0,2) & (1,0) & (1,1) & (1,2) \\
    };
    \node[right=0.3cm of mphy, font=\bfseries] {物理内存排列};

    \draw[brace] (mphy-1-1.south west) -- (mphy-1-3.south east) node[midway, below=5pt, font=\small] {第 0 行};
    \draw[brace] (mphy-1-4.south west) -- (mphy-1-6.south east) node[midway, below=5pt, font=\small] {第 1 行};

    \node[font=\small] (base) at ($(mphy-1-1.south)+(0,-1.3)$) {$\beta$};
    \draw[arrow] (base.north) -- ($(mphy-1-1.south)+(0,-0.05)$);
\end{tikzpicture}
\caption{二维数组行优先 (Row-Major) 压平排列}
\label{fig:array_2d}
\end{figure}

\begin{itemize}
    \item \textbf{稀疏压缩}：$A\in\mathbb{R}^{M\times N}$，三元组 $\mathcal{T}=\{(i,j,A_{ij})\mid A_{ij}\neq0\}$，空间 $O(|\mathcal{T}|)$，检索需索引结构。
    \item \textbf{实例化}：网页 \texttt{sequential-list}、\texttt{matrix} 演示连续寻址与越界；行/列优先、稀疏三元组均为本映射的实例。
\end{itemize}

\subsection{链式存储结构（链表 / Link List）}

\begin{itemize}
    \item \textbf{对象}：节点 $n_k=(d_k,p_k)$，$p_k\in\text{Memory}\cup\{\emptyset\}$。
    \item \textbf{不变式}：$p_k=\text{Addr}(n_{k+1})$，$p_{\text{last}}=\emptyset$；双向加前驱，循环令 $p_{\text{last}}=\text{Addr}(n_1)$。
    \item \textbf{定理}：随机访问 $O(n)$，头插删 $O(1)$。
    \item \textbf{实例化}：网页 \texttt{linked-list}、\texttt{doubly-linked-list}、\texttt{circular-linked-list} 分别演示三类不变式的指针重写。
\end{itemize}

\begin{figure}[H]
\centering
\begin{tikzpicture}[
    head/.style={draw, minimum width=1.1cm, minimum height=0.8cm, font=\ttfamily},
    lnode/.style={rectangle split, rectangle split horizontal, rectangle split parts=2,
                  rectangle split part align={center, center},
                  draw, minimum height=0.8cm, font=\ttfamily, inner sep=6pt},
    arrow/.style={->, >=latex, thick},
    node distance=1.2cm
]
    \node[head] (h) {$h$};
    \node[lnode, right=of h] (a) {$d_1$ \nodepart{two} $\bullet$};
    \node[lnode, right=of a] (b) {$d_2$ \nodepart{two} $\bullet$};
    \node[lnode, right=of b] (c) {$d_3$ \nodepart{two} $\emptyset$};

    \draw[arrow] (h.east) -- (a.west);
    \draw[arrow] (a.east) -- (b.west);
    \draw[arrow] (b.east) -- (c.west);
\end{tikzpicture}
\caption{单向链表 (Singly Linked List)}
\label{fig:linked_singly}
\end{figure}

\begin{figure}[H]
\centering
\begin{tikzpicture}[
    head/.style={draw, minimum width=1.1cm, minimum height=0.8cm, font=\ttfamily},
    dnode/.style={rectangle split, rectangle split horizontal, rectangle split parts=3,
                  rectangle split part align={center, center, center},
                  draw, minimum height=0.8cm, font=\ttfamily, inner sep=5pt},
    arrow/.style={->, >=latex, thick},
    node distance=1.2cm
]
    \node[head] (h) {$h$};
    \node[dnode, right=of h] (a) {$\emptyset$ \nodepart{two} $d_1$ \nodepart{three} $\bullet$};
    \node[dnode, right=of a] (b) {$\bullet$ \nodepart{two} $d_2$ \nodepart{three} $\bullet$};
    \node[dnode, right=of b] (c) {$\bullet$ \nodepart{two} $d_3$ \nodepart{three} $\emptyset$};

    \draw[arrow] (h.east) -- (a.west);
    \draw[arrow] ($(a.east)+(0,0.18)$) -- ($(b.west)+(0,0.18)$);
    \draw[arrow] ($(b.west)-(0,0.18)$) -- ($(a.east)-(0,0.18)$);
    \draw[arrow] ($(b.east)+(0,0.18)$) -- ($(c.west)+(0,0.18)$);
    \draw[arrow] ($(c.west)-(0,0.18)$) -- ($(b.east)-(0,0.18)$);
\end{tikzpicture}
\caption{双向链表 (Doubly Linked List)}
\label{fig:linked_doubly}
\end{figure}

\begin{figure}[H]
\centering
\begin{tikzpicture}[
    head/.style={draw, minimum width=1.1cm, minimum height=0.8cm, font=\ttfamily},
    lnode/.style={rectangle split, rectangle split horizontal, rectangle split parts=2,
                  rectangle split part align={center, center},
                  draw, minimum height=0.8cm, font=\ttfamily, inner sep=6pt},
    arrow/.style={->, >=latex, thick},
    node distance=1.2cm
]
    \node[head] (h) {$h$};
    \node[lnode, right=of h] (a) {$d_1$ \nodepart{two} $\bullet$};
    \node[lnode, right=of a] (b) {$d_2$ \nodepart{two} $\bullet$};
    \node[lnode, right=of b] (c) {$d_3$ \nodepart{two} $\bullet$};

    \draw[arrow] (h.east) -- (a.west);
    \draw[arrow] (a.east) -- (b.west);
    \draw[arrow] (b.east) -- (c.west);
    \coordinate (ret) at ($(h.south)+(0,-0.9)$);
    \draw[arrow, rounded corners=4pt] (c.east) -- ++(0.6,0) |- (ret) -- (h.south);
\end{tikzpicture}
\caption{循环链表 (Circular Linked List)}
\label{fig:linked_circular}
\end{figure}

\section{核心矛盾对比：顺序表 vs 链表}

\begin{table}[htbp]
\centering
\renewcommand{\arraystretch}{1.5}
\begin{tabulary}{\textwidth}{@{}CCC@{}}
\toprule
\textbf{特性维度} & \textbf{顺序存储（数组）} & \textbf{链式存储（链表）} \\
\midrule
\textbf{物理内存要求} & 连续区间 $[\beta,\beta+n\sigma)$ & 离散节点集合 $\{n_1,\dots,n_k\}$ \\
\textbf{元素访问} & $O(1)$：$\text{Addr}(i)=\beta+i\sigma$ & $O(n)$：指针链遍历 \\
\textbf{插入/删除} & $O(n)$：平移 $d_{i+1}\dots d_n$ & $O(1)$：重写常数个指针（已知前驱） \\
\textbf{空间代价} & 连续分配，扩容需 $O(n)$ 复制（摊还） & 额外指针 $O(n\sigma_p)$，但无扩容复制 \\
\bottomrule
\end{tabulary}
\end{table}

\begin{itemize}
    \item \textbf{实例化}：顺序表的连续分配与扩容、链表的指针开销均为分配策略的实例；数学代价不变。
\end{itemize}
